\section{Discussion}

\paragraph{Epoch duration selection.}
Experiment~1 selected a 30-second epoch duration for the Proposed-Med configuration using
macro-F1 on the combined dataset as the primary metric. Importantly, Experiment~3 showed
that macro-F1 was largely insensitive to the epoch grid over 20--120\,s, with only small differences between durations and no evidence of a strong performance cliff. This supports
the practical view that epoch duration is a secondary design choice for the detector
itself; we nevertheless fix 30\,s downstream for consistency and to match the empirically
best setting under the chosen metric.

\paragraph{Strategy comparison.}
Across the five compared conditions, Proposed-Med (macro-F1\,=\,0.8177) did not dominate DBO (macro-F1\,=\,0.8540) under the selected 30-second epoch. Stage~A may remove epochs with atypical but valid blink morphologies, trading recall for robustness. Future work should quantify when Stage~A is beneficial and when it is overly conservative.
Proposed-Med substantially improved over the two practical baselines
(BLINKER-concat, macro-F1\,=\,0.7453; MNE-annot, macro-F1\,=\,0.5610).

\paragraph{Sensitivity to the evaluation strictness.}
Experiment~4 demonstrated that macro-F1 varies by 0.5168 across IoU thresholds $\{0.0,0.1,0.2,0.3,0.5\}$ (stable $<$1\,pp: NO). The main experiments therefore adopt IoU\,=\,0.1 as a pragmatic default,
while reporting the full IoU sweep to make this methodological dependence explicit.

\paragraph{Morphological analysis.}
Experiment~6 (detailed) complements the scalar metrics by visualising typical true
positives, false positives, and false negatives across duration and amplitude strata. This
helps distinguish failures due to boundary misalignment from failures due to missed events,
and provides qualitative evidence about which blink types are systematically challenging.
